\section{Fazit}
Mit dem Raspberry-Pi-Robotik-Controller wurde eine funktionierende Hardwareplattform für Robotikprojekte aufgebaut. Das Grundkonzept funktioniert. Die wichtigsten Bereiche wie Motoren, Servos, Sensoren, Stromversorgung und Kommunikation sind auf einer gemeinsamen Platine zusammengeführt.

Beim Aufbau wurden auch einige Schwachstellen sichtbar. Dazu gehören der Fehler im übernommenen Schaltungsteil, die Empfindlichkeit des 5-V-Wandlers beim Abziehen des Akkus und einzelne offene Punkte im Layout.

Trotzdem ist die aktuelle Platine eine brauchbare Grundlage für die im Antrag vorgesehenen Funktionen. Die Messungen zeigen, welche Bereiche bereits funktionieren und wo noch genauer geprüft werden muss.
